% Pakiety, geometria A5, kroje.
% Przełączniki \ifekran i \ifpolski są ustawiane w main.tex.
\makeatletter

% ---------- Strona (A5, dwustronnie) ----------------------------------
\usepackage[a5paper, inner=18mm, outer=15mm, top=19mm, bottom=20mm,
  headsep=6mm, footskip=10mm, headheight=12.5pt]{geometry}

% ---------- Język -----------------------------------------------------
\ifpolski
  \usepackage[english,polish]{babel}
\else
  \usepackage[polish,english]{babel}
\fi

% ---------- Kroje ------------------------------------------------------
% URW Bookman (klon ITC Bookman) + URW Gothic (klon ITC Avant Garde):
% para, która definiuje typografię lat 70. Obie mają polskie znaki i cyrylicę.
% Pliki .otf: folder fonts/ w katalogu projektu albo zainstalowane w systemie.
% Japoński: fallback na Harano Aji (jest w TeX Live).
\usepackage{fontspec}
\IfFileExists{fonts/URWBookman-Light.otf}
  {\def\bx@fp{fonts/}}{\def\bx@fp{}}

\directlua{
  luaotfload.add_fallback("bxserif", {"[HaranoAjiMincho-Regular.otf]:mode=node;"})
  luaotfload.add_fallback("bxsans",  {"[HaranoAjiGothic-Medium.otf]:mode=node;"})
}

\setmainfont{URWBookman-Light.otf}[
  Path=\bx@fp,
  BoldFont=URWBookman-Demi.otf,
  ItalicFont=URWBookman-LightItalic.otf,
  BoldItalicFont=URWBookman-DemiItalic.otf,
  RawFeature={fallback=bxserif}]
\setsansfont{URWGothic-Book.otf}[
  Path=\bx@fp,
  BoldFont=URWGothic-Demi.otf,
  ItalicFont=URWGothic-BookOblique.otf,
  BoldItalicFont=URWGothic-DemiOblique.otf,
  RawFeature={fallback=bxsans}]
\newfontfamily\bxdisplay{URWGothic-Demi.otf}[
  Path=\bx@fp, RawFeature={fallback=bxsans}]

\usepackage{microtype}

% ---------- Pakiety pomocnicze ----------------------------------------
\usepackage{xcolor}
\usepackage{graphicx}
\graphicspath{{img/}}
\usepackage{tikz}
\usepackage{needspace}
\usepackage{lettrine}
\usepackage{scrlayer-scrpage}
\usepackage{lipsum} % tylko do podglądu — usuń, gdy wstawisz własne teksty
\usepackage[hidelinks, pdfusetitle]{hyperref}

\makeatother
